\section{DISTRIBUTIONS AND TEMPERED DISTRIBUTIONS}

In this section, n denotes a natural integer. We denote by \(\mu\) the Lebesgue measure on \(\mathbf{R}^{n}\) , as well as its restriction to every locally compact subset of \(\mathbf{R}^{n}\) .

We denote

\[
x \cdot y = \sum_ {i = 1} ^ {n} x _ {i} y _ {i}
\]

the scalar product on Euclidean space \(\mathbf{R}^{n}\) ; the Euclidean norm is denoted \(x \mapsto \|x\|\) (TG, VI, p. 7). We recall that the group \(\mathbf{R}^{n}\) is in duality with itself with respect to the map \((x, y) \mapsto \exp(2i\pi x \cdot y)\) and that the dual measure of the Lebesgue measure is then identified with the Lebesgue measure (corollary 3 of II, p. 236). We denote by F (resp. \(\overline{F}\) ) the Fourier transform (resp. the inverse Fourier transform) on \(\mathbf{R}^{n}\) (cf. no. 9 of II, p. 237).

For every \(\alpha = (\alpha_{i})_{1\leqslant i\leqslant n}\in \mathbf{N}^{n}\) , we denote by \(X^{\alpha}\) the function from \(\mathbf{R}^n\) into \(\mathbf{R}\) defined by \(x = (x_{i})_{1\leqslant i\leqslant n}\mapsto x^{\alpha} = \prod_{i = 1}^{n}x_{i}^{\alpha_{i}}\) .

Let \(\alpha\) and \(\beta\) be elements of \(\mathbf{N}^n\) . We denote

\[
| \alpha | = \sum_ {i = 1} ^ {n} \alpha_ {i}, \qquad \binom{\alpha}{\beta} = \prod_ {i = 1} ^ {n} \binom{\alpha_ {i}}{\beta_ {i}}.
\]

Lemma 1. --- Let U be an open subset of \(\mathbf{R}^{n}\) .

\begin{enumerate}
\def\labelenumi{\alph{enumi})}
\item
  Let K be a compact subset of U and V an open neighborhood of K in U. There exists a function \(\varphi \in \mathcal{C}^{\infty}(\mathrm{U})\) with compact support contained in V such that \(0 \leqslant \varphi \leqslant 1\) and such that \(\varphi(x) = 1\) for all \(x \in \mathrm{K}\) ;
\item
  For every locally finite open cover of U, there exists a partition of unity consisting of functions of class \(C^{\infty}\) subordinate to it.
\end{enumerate}

This follows from VAR, R1, p. 40, 5.3.6.

PROPOSITION 1. --- Let U be an open subset of \(\mathbf{R}^{n}\) and \(k\in\mathbf{N}\) . Let E, F, and G be topological vector spaces and b: E×F→G a continuous bilinear map. Let f and g be functions of class \(\mathbf{C}^{k}\) from U into E and F, respectively. Then the map h: x↦b(f(x),g(x)) is of class \(\mathbf{C}^{k}\) on U; moreover, for every \(\alpha\in\mathbf{N}^{n}\) such that \(|\alpha|\leqslant k\) and

for all \(x \in \mathrm{U}\) , we have

\[
\partial^{\alpha}h(x) = \sum_{\substack{\beta \in \mathbf{N}^{n}\\ \beta \leqslant \alpha}}\binom {\alpha}{\beta}b\big(\partial^{\beta}f(x),\partial^{\alpha -\beta}g(x)\big).
\]

The map h is the composition of the map \((f,g)\) : \(U \rightarrow E \times F\) , which is of class \(C^{k}\) , and of the map b: \(E \times F \rightarrow G\) which is of class \(C^{\infty}\) . It is therefore of class \(C^{k}\) . The expression for its partial derivatives results from Leibniz's formula (FVR, I, p. 28, prop. 2, cf. A, III, p. 122, corollary).

We recall that a separated locally convex topological vector space is a Montel space if it is barrelled and if every bounded subset is relatively compact (EVT, IV, p. 18, def. 4).

We also recall that if E and F are locally convex spaces, and if E is bornological (EVT, III, p. 12, def. 1), a linear mapping \(u\colon\mathrm{E}\to\mathrm{F}\) is continuous if and only if the image under \(u\) of every bounded subset of E is bounded in F (EVT, III, p. 11, prop. 1, (iiibis), when F is semi-normed, the general case following formally, cf. EVT, II, p. 7, prop. 5, b)).

